% GENERATED FILE. Edit canonical lessons, then run npm run generate:books.
% canonical-source-hash: fnv1a64:bc4c4f9b37daddca
% canonical-lessons: GU-C337-panu, GU-C337-prakaran, GU-C337-patra, GU-C337-nayak, GU-C337-vacak

\chapter{Turning the Pages}
\label{ch:gu-turning-the-pages}
\hlchaptermodality{337}

\begin{chapteropening}
\textbf{By the end of this chapter:} \emph{I can talk about the parts and people of a book.}

Complete the last lesson of chapter 337: I can talk about the parts and people of a book.
\end{chapteropening}

\section[\texorpdfstring{\gu{પાનું} (pānũ)}{પાનું (pānũ)}]{\gu{પાનું} (pānũ) — a page}

\label{lesson:GU-C337-panu}



\begin{quote}
Before the new one: say the Gujarati for an autobiography, then the Gujarati for imagination.
\end{quote}

\subsection*{You'll want to know: \gu{પાનું}}
\textbf{\gu{પાનું}} — \emph{pānũ} — \textquotedblleft{}a page\textquotedblright{}.

\begin{quote}
\gu{મેં} \gu{પુસ્તકનાં} \gu{સો} \gu{પાનાં} \gu{વાંચ્યાં}.
\end{quote}

\begin{quote}
\emph{mẽ pustaknā̃ so pānā̃ vā̃cyā̃.}
\end{quote}

\begin{quote}
I read a hundred pages of the book.
\end{quote}

A neuter noun.

\begin{cousinweb}
\textbf{\gu{પાનું}} (\emph{pānũ}) goes back to Sanskrit \emph{parṇá}, \textquotedblleft{}a leaf\textquotedblright{}, where CDIAL lists Gujarati \emph{pān}, a leaf, and \emph{pānũ} (CDIAL 7918): a book's leaf is its page. \textbf{\gu{પાન}} is still a leaf on a tree.
\end{cousinweb}

\begin{grammarlens}[title={What you've built}]
One more word for telling what happened.
\end{grammarlens}

\subsection*{Your turn}
\begin{itemize}
\raggedright
  \item \emph{Say it:} \emph{pānũ}
  \item \emph{Say it:} \emph{pānũ}, once more
  \item \emph{Say it:} \emph{mẽ pustaknā̃ so pānā̃ vā̃cyā̃.} — I read a hundred pages of the book.
  \item \emph{Recall:} say the Gujarati for an autobiography, then the Gujarati for imagination, then say \emph{pānũ} again
\end{itemize}

\begin{tcolorbox}[breakable,colback=teal!4,colframe=teal!35!black,title={Before you move on}]
What does \gu{પાનું} mean? (\textquotedblleft{}A page\textquotedblright{}.) Say it once more.
\end{tcolorbox}
\section[\texorpdfstring{\gu{પ્રકરણ} (prakaraṇ)}{પ્રકરણ (prakaraṇ)}]{\gu{પ્રકરણ} (prakaraṇ) — a chapter (of a book)}

\label{lesson:GU-C337-prakaran}



\begin{quote}
Before the new one: say the Gujarati for imagination, then the Gujarati for a page.
\end{quote}

\subsection*{You'll want to know: \gu{પ્રકરણ}}
\textbf{\gu{પ્રકરણ}} — \emph{prakaraṇ} — \textquotedblleft{}a chapter (of a book)\textquotedblright{}.

\begin{quote}
\gu{આ} \gu{પુસ્તકનું} \gu{છેલ્લું} \gu{પ્રકરણ} \gu{સૌથી} \gu{સારું} \gu{છે}.
\end{quote}

\begin{quote}
\emph{ā pustaknũ chellũ prakaraṇ sauthī sārũ che.}
\end{quote}

\begin{quote}
The last chapter of this book is the best.
\end{quote}

A neuter noun.

\begin{cousinweb}
\textbf{\gu{પ્રકરણ}} (\emph{prakaraṇ}) is a tatsama, Sanskrit \emph{prakaraṇa}, \textquotedblleft{}treatment, discussion, a section\textquotedblright{} (Monier-Williams). It is neuter: \textbf{\gu{છેલ્લું} \gu{પ્રકરણ}}.
\end{cousinweb}

\begin{grammarlens}[title={What you've built}]
One more word for telling what happened.
\end{grammarlens}

\subsection*{Your turn}
\begin{itemize}
\raggedright
  \item \emph{Say it:} \emph{prakaraṇ}
  \item \emph{Say it:} \emph{prakaraṇ}, once more
  \item \emph{Say it:} \emph{ā pustaknũ chellũ prakaraṇ sauthī sārũ che.} — The last chapter of this book is the best.
  \item \emph{Recall:} say the Gujarati for imagination, then the Gujarati for a page, then say \emph{prakaraṇ} again
\end{itemize}

\begin{tcolorbox}[breakable,colback=teal!4,colframe=teal!35!black,title={Before you move on}]
What does \gu{પ્રકરણ} mean? (\textquotedblleft{}A chapter (of a book)\textquotedblright{}.) Say it once more.
\end{tcolorbox}
\section[\texorpdfstring{\gu{પાત્ર} (pātra)}{પાત્ર (pātra)}]{\gu{પાત્ર} (pātra) — a character (in a story)}

\label{lesson:GU-C337-patra}



\begin{quote}
Before the new one: say the Gujarati for a page, then the Gujarati for a chapter (of a book).
\end{quote}

\subsection*{You'll want to know: \gu{પાત્ર}}
\textbf{\gu{પાત્ર}} — \emph{pātra} — \textquotedblleft{}a character (in a story)\textquotedblright{}.

\begin{quote}
\gu{આ} \gu{નવલકથામાં} \gu{ઘણાં} \gu{પાત્રો} \gu{છે}.
\end{quote}

\begin{quote}
\emph{ā navalkathāmā̃ ghaṇā̃ pātro che.}
\end{quote}

\begin{quote}
There are many characters in this novel.
\end{quote}

A neuter noun.

\begin{cousinweb}
\textbf{\gu{પાત્ર}} (\emph{pātra}) is a tatsama, Sanskrit \emph{pātra}, \textquotedblleft{}a vessel, a fit receptacle\textquotedblright{} (Monier-Williams); the sense of a person in a story or a play is the modern language's. It is not \textbf{\gu{પત્ર}}, a letter: the vowel is long.
\end{cousinweb}

\begin{grammarlens}[title={What you've built}]
One more word for telling what happened.
\end{grammarlens}

\subsection*{Your turn}
\begin{itemize}
\raggedright
  \item \emph{Say it:} \emph{pātra}
  \item \emph{Say it:} \emph{pātra}, once more
  \item \emph{Say it:} \emph{ā navalkathāmā̃ ghaṇā̃ pātro che.} — There are many characters in this novel.
  \item \emph{Recall:} say the Gujarati for a page, then the Gujarati for a chapter (of a book), then say \emph{pātra} again
\end{itemize}

\begin{tcolorbox}[breakable,colback=teal!4,colframe=teal!35!black,title={Before you move on}]
What does \gu{પાત્ર} mean? (\textquotedblleft{}A character (in a story)\textquotedblright{}.) Say it once more.
\end{tcolorbox}
\section[\texorpdfstring{\gu{નાયક} (nāyak)}{નાયક (nāyak)}]{\gu{નાયક} (nāyak) — the hero (of a story)}

\label{lesson:GU-C337-nayak}



\begin{quote}
Before the new one: say the Gujarati for a chapter (of a book), then the Gujarati for a character.
\end{quote}

\subsection*{You'll want to know: \gu{નાયક}}
\textbf{\gu{નાયક}} — \emph{nāyak} — \textquotedblleft{}the hero (of a story)\textquotedblright{}.

\begin{quote}
\gu{વાર્તાનો} \gu{નાયક} \gu{એક} \gu{ગરીબ} \gu{છોકરો} \gu{છે}.
\end{quote}

\begin{quote}
\emph{vārtāno nāyak ek garīb chokro che.}
\end{quote}

\begin{quote}
The hero of the story is a poor boy.
\end{quote}

A masculine noun.

\begin{cousinweb}
is a tatsama, Sanskrit \emph{nāyaka}, \textquotedblleft{}a leader\textquotedblright{}, and in drama \textquotedblleft{}the hero\textquotedblright{} (Monier-Williams). \textbf{\gu{વીર}} is a brave man; the \textbf{\gu{નાયક}} is the main person of a story.
\end{cousinweb}

\begin{grammarlens}[title={What you've built}]
One more word for telling what happened.
\end{grammarlens}

\subsection*{Your turn}
\begin{itemize}
\raggedright
  \item \emph{Say it:} \emph{nāyak}
  \item \emph{Say it:} \emph{nāyak}, once more
  \item \emph{Say it:} \emph{vārtāno nāyak ek garīb chokro che.} — The hero of the story is a poor boy.
  \item \emph{Recall:} say the Gujarati for a chapter (of a book), then the Gujarati for a character, then say \emph{nāyak} again
\end{itemize}

\begin{tcolorbox}[breakable,colback=teal!4,colframe=teal!35!black,title={Before you move on}]
What does \gu{નાયક} mean? (\textquotedblleft{}The hero (of a story)\textquotedblright{}.) Say it once more.
\end{tcolorbox}
\section[\texorpdfstring{\gu{વાચક} (vācak)}{વાચક (vācak)}]{\gu{વાચક} (vācak) — a reader}

\label{lesson:GU-C337-vacak}



\begin{quote}
Before the new one: say the Gujarati for a character, then the Gujarati for the hero (of a story).
\end{quote}

\subsection*{You'll want to know: \gu{વાચક}}
\textbf{\gu{વાચક}} — \emph{vācak} — \textquotedblleft{}a reader\textquotedblright{}.

\begin{quote}
\gu{આ} \gu{પુસ્તક} \gu{વાચકોને} \gu{ખૂબ} \gu{ગમ્યું}.
\end{quote}

\begin{quote}
\emph{ā pustak vāckone khūb gamyũ.}
\end{quote}

\begin{quote}
Readers loved this book.
\end{quote}

A masculine or feminine noun, by the person meant.

\begin{cousinweb}
\textbf{\gu{વાચક}} (\emph{vācak}) is a tatsama, Sanskrit \emph{vācaka}, \textquotedblleft{}speaking, telling\textquotedblright{} (Monier-Williams); in Gujarati it is the one who reads, close to \textbf{\gu{વાંચવું}}, to read.
\end{cousinweb}

\begin{grammarlens}[title={What you've built}]
One more word for telling what happened.
\end{grammarlens}

\subsection*{Your turn}
\begin{itemize}
\raggedright
  \item \emph{Say it:} \emph{vācak}
  \item \emph{Say it:} \emph{vācak}, once more
  \item \emph{Say it:} \emph{ā pustak vāckone khūb gamyũ.} — Readers loved this book.
  \item \emph{Recall:} say the Gujarati for a character, then the Gujarati for the hero (of a story), then say \emph{vācak} again
\end{itemize}

\begin{tcolorbox}[breakable,colback=teal!4,colframe=teal!35!black,title={Before you move on}]
What does \gu{વાચક} mean? (\textquotedblleft{}A reader\textquotedblright{}.) Say it once more.
\end{tcolorbox}
